\documentclass[a4paper]{article}
\usepackage[margin=.8in]{geometry}
\usepackage{amsmath}
\usepackage{enumitem}
\usepackage[cache=true]{minted}
\usepackage{amsfonts}
\usepackage{amssymb}
\usepackage{mathtools}
\usepackage{amsmath}
\input{../../macros}
\usepackage{fancyhdr}
\pagestyle{fancy}
\fancyhf{}
\lhead{Name: }
\chead{Quiz 3 -- Interpolation and Approximation (practice quiz)}
\rhead{\today}

\begin{document}
\begin{enumerate}

    \item
        Suppose that $x_0 < x_1 < \dotsb < x_n$ and $y_0, y_1, \dotsc, y_n$ are given real numbers.
        Then there exists a unique polynomial $p \in \mathcal P_n$ such that $p(x_i) = y_i$ for all $i \in \{0, \dotsc, n\}$.

    \item
        Suppose that $f \colon [-1, 1] \to \mathbb R$ is given by $f(x) = \e^{-2x}$,
        and for any $n \in \mathbb N$,
        let $f_n \in \mathcal P_n$ denote the polynomial interpolating~$f$ at $n+1$ equidistant points $-1 = x_0 < x_1 < \dotsc < x_n = 1$.
        Then
        \[
            \lim_{n \to \infty} \left( \max_{-1 \leq x \leq 1} \bigl\lvert f(x) - f_n(x) \bigr\rvert \right) = 0.
        \]

    \item
        Define $\mathcal F \colon \mathcal P_n \to \real$ as the function which,
        to a polynomial $p \in \mathcal P_n$, associates its maximum absolute value over the interval $[-1, 1]$,
        that is
        \[
            \mathcal F(p) = \max_{x \in [-1, 1]} \left\lvert p(x) \right\rvert \, .
        \]
        Then it holds that $\mathcal F(p) \geq \mathcal F(T_n)$ for all $p \in \mathcal P_n$,
        where $T_n$ is the Chebyshev polynomial of degree~$n$.


    \item
        Suppose that $f \colon [-1, 1] \to \mathbb R$ is the Runge function $f(x) = \frac{1}{1 + 25 x^2}$.
        For any $n \in \mathbb N$,
        let $f_n \in \mathcal P_n$ denote the polynomial interpolating~$f$ at $n+1$ \emph{equidistant} points,
        and let $g_n \in \mathcal P_n$ denote the polynomial interpolating~$f$ at the $n+1$ \emph{Chebyshev} points (the roots of $T_{n+1}$).
        Then it holds that
        \[
            \lim_{n \to \infty} \left( \max_{-1 \leq x \leq 1} \bigl\lvert f(x) - f_n(x) \bigr\rvert \right) = +\infty \, ,
            \qquad
            \lim_{n \to \infty} \left( \max_{-1 \leq x \leq 1} \bigl\lvert f(x) - g_n(x) \bigr\rvert \right) = 0 \, .
        \]

    \item
        Suppose that $x_0 < x_1 < \dotsb < x_n$ and $y_0, y_1, \dotsc, y_n$ are given real numbers.
        Then for any $m \in \{0, 1, \dotsc\}$,
        there exists exactly one polynomial $p \in \mathcal P_m$ such that 
        \[
            \frac{1}{n} \sum_{i=0}^{n} \Bigl\lvert p(x_i) - y_i \Bigr\rvert^2 = 0.
        \]

    \item
        Given $x_0 < x_1 < \dotsc < x_n \in \mathbb R$
        and $y_0, y_1, \dotsc , y_n \in \mathbb R$,
        and let $m \geq n$. 
        There exists a unique $p \in \mathcal P_m$ such that the following quantity is minimized:
        \[
            \frac{1}{n} \sum_{i=0}^{n} \Bigl\lvert p(x_i) - y_i \Bigr\rvert^2.
        \]

    \item
        Given $x_0 < x_1 < \dotsc < x_n \in \mathbb R$
        and $y_0, y_1, \dotsc , y_n \in \mathbb R$ with $n = 10$,
        there exists a unique polynomial $p(x) = ax + b$ that minimizes
        \[
            J(a, b) = \frac{1}{2} \sum_{i=0}^n \bigl\lvert y_i - p(x_i) \bigr\rvert^2.
        \]

    \item
        In Julia, if \julia{v} is a vector of size 5,
        then \julia{v[[1, 2, 5]]} returns a vector with the first, second and fifth element of \julia{v}.

    \item
        In Julia, if \julia{v} is a vector of size 5,
        then \julia{v[[true, true, false, true, false]]} returns a vector with the first, second and fourth element of \julia{v}.

    \item
        In Julia, if \julia{A} is a $10\times10$ matrix,
        then \julia{A[mod.(1:end, 2) .== 0, mod.(1:end, 2) .== 0]} gives the matrix obtained by removing the second column and the second row.
\end{enumerate}

\ifdefined\showsolutions
\newpage
\lhead{Instructor copy}
\section*{Answer key}
\begin{enumerate}[itemsep=8pt]
    \item \textbf{True.}
        This is the fundamental theorem of interpolation: the set of nodes
        $x_0 < x_1 < \dotsb < x_n$ leads to a nonsingular Vandermonde
        system, so there is a unique $p \in \mathcal P_n$ solving
        $p(x_i) = y_i$ for all $i$.

    \item \textbf{True.}
        The function $x \mapsto \e^{-2x}$ is entire, so polynomial
        interpolation on equidistant nodes converges uniformly on $[-1,1]$
        as $n \to \infty$.

    \item \textbf{False.}
        The minimal property goes in the \emph{opposite} direction:
        among monic polynomials of degree $n$, $T_n$ has the smallest
        possible sup-norm on $[-1,1]$.
        For example, the zero polynomial satisfies
        $\mathcal F(0) = 0 < 1 = \mathcal F(T_n)$.

    \item \textbf{True.}
        This is the statement of the Runge phenomenon: interpolation at
        equidistant points diverges for the Runge function, while
        interpolation at the Chebyshev nodes (the roots of $T_{n+1}$)
        converges uniformly.

    \item \textbf{False.}
        For $m < n$ interpolation with a degree-$m$ polynomial is in
        general impossible (the condition imposes $n+1$ equalities), and
        for $m > n$ the polynomial is not unique.

    \item \textbf{False.}
        When $m > n$ there are more coefficients than data points, so the
        design matrix does not have full column rank and the
        least-squares problem has infinitely many minimizers
        (indeed the minimum can be zero, attained by the interpolating
        polynomial).

    \item \textbf{True.}
        With 11 data points and only 2 unknown parameters $a, b$, the
        normal equations are a $2 \times 2$ system whose matrix is
        nonsingular as long as the nodes are not all identical, so the
        minimizer is unique.

    \item \textbf{True.}
        Indexing a vector with a vector of indices selects the
        corresponding entries in that order: here the first, second and
        fifth elements of \julia{v}.

    \item \textbf{True.}
        Indexing with a Boolean mask of the same length keeps exactly the
        entries in the positions where the mask is \julia{true}: the
        first, second and fourth elements.

    \item \textbf{False.}
        \julia{mod.(1:end, 2) .== 0} produces the mask
        \julia{[false, true, false, true, ...]}, which selects the
        even-indexed rows and the even-indexed columns, giving the $5 \times 5$
        submatrix of rows and columns $2, 4, 6, 8, 10$ --- not the matrix
        obtained by deleting row 2 and column 2.
\end{enumerate}
\fi

\end{document}
